\chapter{Rotations, Spinors, and the Road to Color}

We begin with rotations because where we are going is color. Leonard Susskind chooses the most concrete possible road into the subject: start with a symmetry we can picture, let it act on objects we already know, and then watch the same mathematics reappear in less familiar settings. These notes follow that lecture closely, with chapter curation by LazyingArt LLC. The thread running through the whole discussion is simple and powerful: the same group-theoretic structure can act on spatial vectors, on spin states, and finally on quark color.

\section{Rotations as the Warm-Up for Color}

Rotations of space are symmetries, and one of the first things they act on is an ordinary vector. That is the point of departure. But it is not the destination. Susskind says early on where he is heading: to color, meaning the red, green, and blue states of a quark. Color is not a two-state system like spin up and spin down. It is a three-state system, and the interesting operations are the ones that mix those three states among themselves.

So the lecture does not study rotations for their own sake. It studies rotations because they give us a familiar example of a symmetry group together with its representations. Once that pattern is in hand, the same logic can be lifted from space to spin and from spin to color.

\section{Axis, Angle, and the Group of Rotations}

How do we parameterize a rotation? Susskind begins in the most geometric language. A rotation has an axis, which we describe by a unit vector \(\hat n\), and it has an angle \(\theta\) through which we rotate about that axis. The unit vector requires two parameters, which one may picture as longitude and latitude on the sphere, and the angle gives one more. Thus a rotation is a three-parameter object.

\begin{figure}[htbp]
\centering
\includegraphics[width=0.58\textwidth]{lecture_03_figure_02.png}
\caption{Unit-vector axis sketch for a rotation.}
\end{figure}

The board sketch is deliberately minimal. It shows the three spatial axes, a distinguished direction \(\hat n\), and a nearby \(\theta\). The lecture speaks of a unit vector pointing to a unit sphere, but the frame itself shows only the axis sketch and the direction arrow.

\begin{figure}[htbp]
\centering
\begin{tikzpicture}[scale=1.2]
\draw[thick] (0,0) -- (0,2.0);
\draw[thick] (0,0) -- (2.1,0);
\draw[thick] (0,0) -- (-1.2,-1.6);
\draw[->,thick] (0,0) -- (1.35,1.15);
\node at (1.55,1.28) {$\hat n$};
\node at (1.95,1.18) {$\theta$};
\end{tikzpicture}
\caption{Cleaned reconstruction of the axis-angle sketch.}
\end{figure}

The lecture also fixes a convention: we take rotations to be right-handed. A left-handed rotation may be rephrased either as a right-handed rotation by a negative angle or as a right-handed rotation about the opposite axis. The point here is not to settle the global topology of the rotation group. It is only to isolate the local data \((\hat n,\theta)\) that carry the three parameters.

\subsection{Question \& Answer}

\paragraph{Question.} Do rotations really form a group, and what exactly makes them non-abelian?

\paragraph{Answer.} Yes. If we rotate an object about one axis and then about another, the result is again some rotation. That gives closure. There is an identity element, namely no rotation at all. Every rotation has an inverse, obtained by rotating through the opposite angle about the same axis. Associativity also holds.

What fails is commutativity. A rotation about the \(x\)-axis followed by a rotation about the \(y\)-axis is not, in general, the same as the same two rotations performed in the opposite order. These are sequential operations, one after another. The lecture even remarks that for very small rotations the failure of commutativity becomes less dramatic, but the group remains unmistakably non-abelian. This is the first non-abelian group encountered in the course.

\section{Vectors and the \(3\times 3\) Representation}

Once the group structure is on the table, Susskind asks for a representation. The simplest one is the representation on the components of an ordinary spatial vector. Let \(V_i\) be the components of a vector \(V\), with \(i=1,2,3\) or, if we prefer, \(i=x,y,z\). For every rotation there is a corresponding \(3\times 3\) matrix \(R_{ij}(\theta,\hat n)\) such that
\begin{equation}
V_i' = R_{ij}(\theta,\hat n)\,V_j.
\end{equation}

\begin{figure}[htbp]
\centering
\includegraphics[width=0.68\textwidth]{lecture_03_figure_03.png}
\caption{Rotation matrix acting on vector components.}
\end{figure}

This is the bridge from geometry to algebra. The board writes the equation in the informal order
\[
R_{ij}(\theta,\hat n)\,V_j = V_i',
\]
and we standardize it as
\[
V_i' = R_{ij}(\theta,\hat n)\,V_j.
\]
For the moment the vector components and the matrix entries are taken to be real.

Now let us do what the lecture does on the board: let us see what invariant length says. A true rotation does not change the length of a vector, so we require
\begin{equation}
V_iV_i = V_i'V_i'.
\end{equation}
Substituting the transformation law gives
\begin{align}
V_i'V_i'
  &= \bigl(R_{ij}V_j\bigr)\bigl(R_{ik}V_k\bigr) \\
  &= R_{ij}R_{ik}V_jV_k .
\end{align}
This must agree with \(V_iV_i\) for every vector \(V\). Therefore the coefficient multiplying \(V_jV_k\) must be the Kronecker delta:
\begin{equation}
R_{ij}R_{ik}=\delta_{jk}.
\end{equation}
In matrix language this becomes
\begin{equation}
R^T R = \mathbf 1.
\end{equation}
That is the defining property of a rotation matrix. It follows immediately that
\begin{equation}
R^{-1}=R^T.
\end{equation}

This is the lecture's first real derivation, and it sets the pattern for everything that follows. We first ask what the transformation is supposed to preserve; then we translate that preservation law into a condition on the matrix representing the transformation.

It is also worth noting the broader point Susskind makes here. This \(3\times 3\) representation is not the only representation of the rotation group. It is simply the one that is easiest to visualize.

\section{Why Spin-1 Particles Are Vector Particles}

With the vector representation in hand, the lecture makes an important conceptual jump. A spin-one particle has three spin states. Therefore the state of the spin alone may be described by a three-component column vector, whose entries are the amplitudes for those three states. If we rotate the spin-one system, those three amplitudes must mix among themselves under a three-dimensional representation of the rotation group.

And in fact it is just this one. The same three-dimensional mathematics that rotates the components of a spatial vector rotates the three spin states of a spin-one particle. That is the content of the phrase that spin-one particles are vector particles. It does not mean that the state is literally a geometric arrow. It means that the three states transform the same way the three components of a vector do.

At this point the lecture pauses and takes a breath. We have already met two different kinds of groups. Some are discrete: one can either do the operation or not do it, as with a reflection. Others are continuous: the group elements depend on continuously varying parameters, as with rotations or multiplication by a phase \(e^{i\phi}\). The continuous groups are the ones that will dominate the rest of the course, and they are called Lie groups.

This pause matters, because it resets the discussion just before the harder case. One can also insert here the trivial case of spin zero. A spin-zero particle is unchanged by a rotation. Its representation of the rotation group is just the identity operator and nothing more.

\section{Spin-$\frac12$ States and Unitary Matrices}

Now we turn to the genuinely new case: spin \(\frac12\). A spin-\(\frac12\) system has two basis states,
\begin{equation}
\begin{pmatrix}1\\0\end{pmatrix},
\qquad
\begin{pmatrix}0\\1\end{pmatrix},
\end{equation}
which we identify with spin up and spin down relative to some chosen axis, say the \(z\)-axis. A general state is
\begin{equation}
\psi =
\begin{pmatrix}
\alpha_1\\
\alpha_2
\end{pmatrix},
\qquad
|\alpha_1|^2+|\alpha_2|^2=1.
\end{equation}
The two entries are complex. So at first sight there are four real parameters. But one overall phase is physically irrelevant, and the normalization condition removes one more. That leaves two physically meaningful parameters, just enough to specify a direction. The lecture does not stop to develop that mapping explicitly; it simply uses the fact that this two-component column is the object that rotations act on.

So what does it mean to rotate such a state? The answer is the same in form as before: we act with a matrix. There is a \(2\times 2\) matrix \(U\) such that
\begin{equation}
\alpha_i' = U_{ij}\alpha_j.
\end{equation}
The logic again begins with preservation, but now what must be preserved is not ordinary Euclidean length. It is the total probability. A state must go into another state.

Let us therefore work out the analogue of the vector derivation. We require
\begin{equation}
\alpha_i^*\alpha_i = \alpha_i'^*\alpha_i'.
\end{equation}
Using \(\alpha_i' = U_{ij}\alpha_j\), we find
\begin{align}
\alpha_i'^*\alpha_i'
  &= \bigl(U_{ij}\alpha_j\bigr)^*\bigl(U_{ik}\alpha_k\bigr) \\
  &= \alpha_j^* U_{ij}^* U_{ik}\alpha_k \\
  &= \alpha_j^*\bigl[(U^T)^*U\bigr]_{jk}\alpha_k .
\end{align}
If this is to agree with \(\alpha_i^*\alpha_i\) for every state, then the matrix multiplying \(\alpha_j^*\alpha_k\) must be the identity:
\begin{equation}
\left(U^T\right)^*U = \mathbf 1.
\end{equation}

\begin{figure}[htbp]
\centering
\includegraphics[width=0.62\textwidth]{lecture_03_figure_04.png}
\caption{Unitary condition written in Hermitian-conjugate form.}
\end{figure}

This is exactly what appears on the board. The secure blackboard content is the pair of equations
\begin{equation}
\left(U^T\right)^*U = \mathbf 1,
\qquad
U^\dagger U = \mathbf 1.
\end{equation}
The second line is simply the compact rewrite of the first, once we name the transpose-plus-complex-conjugate operation:
\begin{equation}
U^\dagger = (U^T)^*.
\end{equation}
Thus the matrices that rotate a spin-\(\frac12\) state are unitary matrices.

The pattern is the same as before, but the setting is different. For vectors we preserved \(V_iV_i\). For spinors we preserve \(\alpha_i^*\alpha_i\). That single difference is exactly what replaces transpose by Hermitian conjugation.

\section{From \(U(2)\) to \(SU(2)\), and Where the Sigmas Enter}

The lecture now performs one of its characteristic parameter counts. A \(2\times 2\) complex matrix has four complex entries, hence eight real parameters. The condition
\[
U^\dagger U=\mathbf 1
\]
is a matrix equation, so it gives four real constraints. Therefore a general unitary \(2\times 2\) matrix has four real parameters.

But a rotation has only three parameters. So there is one parameter too many. That is the problem the lecture wants us to feel before it introduces the next restriction.

What can that extra parameter be? Unitarity tells us only that
\begin{equation}
\det U\,\det U^\dagger = 1,
\end{equation}
and since \(\det U^\dagger = (\det U)^*\), it follows that \(|\det U|=1\). In other words, before any further restriction,
\begin{equation}
\det U = e^{i\phi}
\end{equation}
for some phase \(\phi\). The lecture then fixes that phase by imposing
\begin{equation}
\det U = 1.
\end{equation}
This removes exactly one real parameter.

Why is this consistent with the group structure? Because determinants multiply:
\begin{equation}
\det(AB)=\det A\,\det B.
\end{equation}
So if two matrices both have determinant \(1\), their product also has determinant \(1\). The identity matrix has determinant \(1\) as well. Thus the unitary matrices with determinant \(1\) form a closed subgroup.

This is the special unitary group. In the notation used throughout physics,
\begin{equation}
U(n): \text{unitary } n\times n \text{ matrices}, \qquad
SU(n): \text{unitary } n\times n \text{ matrices with } \det U=1.
\end{equation}
A \(1\times 1\) unitary matrix is just a phase \(e^{i\phi}\), the group \(U(1)\). The group acting on the two components of a spinor is \(SU(2)\).

At lecture level Susskind then says something deliberately forceful: this is the rotation group. Immediately after saying it, however, he adds that there is a subtlety he is not discussing. The safe way to record the point is this: \(SU(2)\) has the correct three local parameters and the correct multiplication structure for describing spatial rotations on spin-\(\frac12\) states, even though it does not act on ordinary three-vectors.

\subsection{Question \& Answer}

\paragraph{Question.} Where do the Pauli matrices come in?

\paragraph{Answer.} They come in when we stop looking at finite transformations as whole group elements and instead examine transformations very close to the identity. Small rotations should be represented by matrices very close to \(\mathbf 1\). So let us write
\begin{equation}
U=\mathbf 1+\epsilon M+O(\epsilon^2),
\end{equation}
with \(\epsilon\) small and real.

Now impose unitarity to first order. Since
\[
U^\dagger = \mathbf 1+\epsilon M^\dagger+O(\epsilon^2),
\]
we have
\begin{equation}
U^\dagger U
 = \mathbf 1+\epsilon(M^\dagger+M)+O(\epsilon^2).
\end{equation}
For this to equal \(\mathbf 1\), we must have
\begin{equation}
M^\dagger=-M.
\end{equation}
So in this convention the infinitesimal generator is anti-Hermitian.

Susskind then makes the standard improvement and pulls out a factor of \(i\). Write instead
\begin{equation}
U=\mathbf 1+i\epsilon M+O(\epsilon^2).
\end{equation}
Then
\[
U^\dagger = \mathbf 1-i\epsilon M^\dagger+O(\epsilon^2),
\]
and therefore
\begin{equation}
U^\dagger U
 = \mathbf 1+i\epsilon(M-M^\dagger)+O(\epsilon^2).
\end{equation}
Thus first-order unitarity now requires
\begin{equation}
M^\dagger=M.
\end{equation}
So after pulling out the factor of \(i\), the generator becomes Hermitian.

We still have to impose determinant one. Near the identity, the lecture analyzes the determinant by a direct \(2\times 2\) expansion. The clean first-order statement is
\begin{equation}
\det(\mathbf 1+\epsilon M)=1+\epsilon\,\operatorname{tr}M+O(\epsilon^2).
\end{equation}
Accordingly,
\begin{equation}
\det(\mathbf 1+i\epsilon M)=1+i\epsilon\,\operatorname{tr}M+O(\epsilon^2),
\end{equation}
and the condition \(\det U=1\) implies
\begin{equation}
\operatorname{tr}M=0.
\end{equation}

So the generators of small \(SU(2)\) rotations are precisely the traceless Hermitian \(2\times 2\) matrices. That is exactly where the sigmas enter. The Pauli matrices span this space, so every such generator can be written as
\begin{equation}
M=a_x\sigma_x+a_y\sigma_y+a_z\sigma_z.
\end{equation}
If the earlier spin matrices are denoted \(S_i\), then they are proportional to these sigmas; at lecture level the only point needed here is that the sigmas furnish the basis of the infinitesimal generators. A small rotation about the \(x\)-axis uses \(\sigma_x\), about the \(y\)-axis uses \(\sigma_y\), and a rotation about a general axis uses the corresponding linear combination.

This is the algebraic heart of the matter. The Pauli matrices are not decorative extra objects added from elsewhere. They appear because \(SU(2)\) near the identity demands traceless Hermitian \(2\times 2\) generators, and that is exactly what the sigmas are.

\section{Color as a Three-State System and the Arrival of \(SU(3)\)}

Only after all of this preparation does the lecture return to its announced target. A quark has three color states: red, green, and blue. So the color state of a single quark is naturally described by a three-component column,
\begin{equation}
q=
\begin{pmatrix}
\alpha_1\\
\alpha_2\\
\alpha_3
\end{pmatrix}
\qquad \text{or simply} \qquad
q=
\begin{pmatrix}
r\\
g\\
b
\end{pmatrix}.
\end{equation}
The entries may be read as amplitudes, so that the squared magnitudes give the probabilities for the quark to be red, green, or blue. Susskind also remarks that one may think similarly of three field operators, one for each color, assembled into a three-component object.

One can certainly imagine discrete symmetries that interchange two colors. But the lecture wants the larger continuous family that contains such mixing in a more powerful way. The natural generalization of the spinor story is to let special unitary \(3\times 3\) matrices act on the color triplet:
\begin{equation}
q' = Uq,
\qquad
U^\dagger U=\mathbf 1,
\qquad
\det U=1,
\qquad
U\in SU(3).
\end{equation}
This is the same mathematics as before, except that the matrices are now \(3\times 3\) rather than \(2\times 2\). It is not a rotation in ordinary space. It is a rotation in an internal three-state space.

The lecture then counts parameters just as it did for \(SU(2)\). A \(3\times 3\) complex matrix has nine complex entries, hence eighteen real parameters. The condition
\[
U^\dagger U=\mathbf 1
\]
gives nine real constraints at lecture level, one for each matrix entry, and the condition \(\det U=1\) removes one more. So the group has
\begin{equation}
18-9-1=8
\end{equation}
independent parameters.

This is the first moment at which the number eight appears in a serious way. Susskind immediately points out the suggestive connection: if one is wondering what this eight has to do with the eight gluons, one is wondering about the right thing. But the lecture stops short of giving the full gauge-theory construction. Even the joke that there are not nine gluons because ``we threw away the determinant'' is meant heuristically, as a pointer toward the eight-generator structure of \(SU(3)\), not as a finished derivation.

What matters for the chapter at this stage is the symmetry statement itself. The claim is that physics is invariant under these \(SU(3)\) transformations of color. That is the symmetry group of quantum chromodynamics. Just as the earlier parts of the course used symmetry to constrain what could appear in a theory, the next step will be to demand that the QCD Lagrangian be invariant under this \(SU(3)\) action.

\section{Summary}

The lecture unfolds in a remarkably controlled sequence. We begin with ordinary rotations, parameterized by an axis \(\hat n\) and an angle \(\theta\), and we verify that they form a three-parameter non-abelian group. We then represent them on vector components and derive the defining condition
\[
R^T R=\mathbf 1
\]
from invariant length. That same three-dimensional representation reappears as the action of rotations on the three states of a spin-one particle.

After a brief pause to distinguish discrete from continuous symmetries, the lecture turns to spin \(\frac12\). There the state is a two-component complex spinor, and the preserved quantity is total probability rather than Euclidean length. The corresponding matrix condition is
\[
U^\dagger U=\mathbf 1.
\]
Parameter counting then shows that one further restriction is needed, and \(\det U=1\) leads to \(SU(2)\). Near the identity, the generators become traceless Hermitian \(2\times 2\) matrices, hence linear combinations of the Pauli matrices.

Only then do we arrive at color. A quark color state is a three-component object, and the same special-unitary logic lifts from \(SU(2)\) to \(SU(3)\). The resulting eight-parameter symmetry is the first mathematical sign of the gluonic structure still to come. The lecture ends exactly there: not with the full construction of QCD, but with the symmetry structure now plainly visible.